\clearpage
\begin{trapbox}
\textbf{\textcolor{CTrap}{容易出错的地方，不只要记“不能”，还要知道为什么}}
\textbf{\textcolor{CTrap}{易错1：误以为每个二次函数都能写成实数两根式。}}
\begin{center}\begin{tikzpicture}[font=\small,>=stealth,line join=round]
\draw[gray!40,rounded corners=3pt] (0,0) rectangle (14.2,8.5);\node at (7.1,7.9) {开口向上：顶点高低决定实数交点个数};
\node at (2.4,6.9) {$\Delta=4>0$};\begin{scope}[shift={(2.4,3.8)},x=0.85cm,y=0.6cm]\draw[->,gray] (-2,0)--(2,0) node[right]{$x$};\draw[->,gray] (0,-2)--(0,4) node[above]{$y$};\draw[blue,thick,solid,domain=-1.7:1.7,samples=160] plot(\x,{(\x)^2-1});\fill[orange!85!black] (-1,0) circle(2pt);\fill[orange!85!black] (1,0) circle(2pt);\fill[orange!85!black] (0,-1) circle(2pt);\end{scope}\node[align=center,text width=4cm] at (2.4,1.7) {$y=x^2-1$\\两个交点};\node at (7.050000000000001,6.9) {$\Delta=0$};\begin{scope}[shift={(7.050000000000001,3.8)},x=0.85cm,y=0.6cm]\draw[->,gray] (-2,0)--(2,0) node[right]{$x$};\draw[->,gray] (0,-2)--(0,4) node[above]{$y$};\draw[teal,thick,solid,domain=-1.7:1.7,samples=160] plot(\x,{(\x)^2});\fill[orange!85!black] (0,0) circle(2pt);\end{scope}\node[align=center,text width=4cm] at (7.050000000000001,1.7) {$y=x^2$\\一个交点};\node at (11.7,6.9) {$\Delta=-4<0$};\begin{scope}[shift={(11.7,3.8)},x=0.85cm,y=0.6cm]\draw[->,gray] (-2,0)--(2,0) node[right]{$x$};\draw[->,gray] (0,-2)--(0,4) node[above]{$y$};\draw[violet,thick,solid,domain=-1.7:1.7,samples=160] plot(\x,{(\x)^2+1});\fill[orange!85!black] (0,1) circle(2pt);\end{scope}\node[align=center,text width=4cm] at (11.7,1.7) {$y=x^2+1$\\没有交点};
\node[text width=12.6cm,align=center] at (7.1,.65) {同尺度，$a=1$：$k=-1,0,1$，恒有$\Delta=-4ak$。};
\end{tikzpicture}\end{center}


例如 $x^2+1>0$，从不等于 $0$，所以没有实数根；正确做法是先判断 $\Delta\ge0$ 是否成立。


\end{trapbox}
\clearpage
\begin{trapbox}
\textbf{\textcolor{CTrap}{易错2：误以为 $\Delta=0$ 时有两个不同交点。}}
\begin{center}\begin{tikzpicture}[font=\small,>=stealth,line join=round]
\draw[gray!40,rounded corners=3pt] (0,0) rectangle (14.2,8.3);\node at (7.1,7.700000000000001) {重根：两个括号相同，图上只有一个交点};
\begin{scope}[shift={(1.7000000000000002,2.57)},x=1.2cm,y=0.9199999999999999cm]\draw[->,gray] (-0.5,0)--(4.5,0) node[right]{$x$};\draw[->,gray] (0,-1)--(0,4) node[above]{$y$};\draw[blue,thick,solid,domain=0.2:3.8,samples=180] plot(\x,{(\x-2)^2});\fill[orange!85!black] (2,0) circle(2.2pt);\node[below=4pt] at (2,0){$(2,0)$};\end{scope}\node[anchor=west,align=left,text width=5.0cm] at (8.4,4.0) {$f=(x-2)^2$\\[9pt]$=(x-2)(x-2)$\\[9pt]$x_1=x_2=2$\\[9pt]$\Delta=0$\\[9pt]交点只有 $(2,0)$};
\node[text width=12.6cm,align=center] at (7.1,.65) {两个相同的根不能计为两个不同的交点。};
\end{tikzpicture}\end{center}


例如 $(x-2)^2$ 只有一个交点 $(2,0)$；写成两根式时，是 $(x-2)(x-2)$，两个根相同。


\end{trapbox}
\clearpage
\begin{trapbox}
\textbf{\textcolor{CTrap}{易错3：写错对称轴的正负号。}}
\begin{center}\begin{tikzpicture}[font=\small,>=stealth,line join=round]
\draw[gray!40,rounded corners=3pt] (0,0) rectangle (14.2,8.3);\node at (7.1,7.700000000000001) {三种写法，对应同一条抛物线};
\begin{scope}[shift={(1.7000000000000002,2.9642857142857144)},x=1.2cm,y=0.6571428571428571cm]\draw[->,gray] (-0.5,0)--(4.5,0) node[right]{$x$};\draw[->,gray] (0,-2)--(0,5) node[above]{$y$};\draw[blue,thick,solid,domain=-0.2:4.2,samples=180] plot(\x,{(\x-2)^2-1});\fill[orange!85!black] (0,3) circle(2.2pt);\fill[orange!85!black] (1,0) circle(2.2pt);\fill[orange!85!black] (3,0) circle(2.2pt);\fill[orange!85!black] (2,-1) circle(2.2pt);\draw[violet,dashed] (2,-.95)--(2,4.5);\node[left] at (0,3){$(0,3)$};\node[below=5pt] at (1,0){$1$};\node[below=5pt] at (3,0){$3$};\node[below=4pt] at (2,-1){$(2,-1)$};\end{scope}\node[anchor=west,align=left,text width=5.0cm] at (8.4,4.0) {$x^2-4x+3$\\[9pt]$=(x-2)^2-1$\\[9pt]$=(x-1)(x-3)$\\[9pt]顶点 $(2,-1)$\\[9pt]交点 $(1,0),(3,0)$};
\node[text width=12.6cm,align=center] at (7.1,.65) {一般式读纵轴交点，顶点式读顶点，两根式读横轴交点。};
\end{tikzpicture}\end{center}


对于 $x^2-4x+3$，$h=-(-4)/(2\cdot1)=2$，不是 $-2$。不放心就配方成 $(x-2)^2-1$ 检查。


\end{trapbox}
\clearpage
\begin{trapbox}
\textbf{\textcolor{CTrap}{易错4：误以为顶点式里的 $k$ 一定是最小值。}若 $y=-(x-1)^2+3$，开口向下，$3$ 是\textbf{最大值}}
\begin{center}\begin{tikzpicture}[font=\small,>=stealth,line join=round]
\draw[gray!40,rounded corners=3pt] (0,0) rectangle (14.2,8.3);\node at (7.1,7.700000000000001) {顶点式中的k：开口向下时是最大值};
\begin{scope}[shift={(3.0285714285714285,3.375)},x=1.0714285714285714cm,y=0.575cm]\draw[->,gray] (-1.8,0)--(3.8,0) node[right]{$x$};\draw[->,gray] (0,-3)--(0,5) node[above]{$y$};\draw[blue,thick,solid,domain=-1.3:3.3,samples=180] plot(\x,{-(\x-1)^2+3});\fill[orange!85!black] (1,3) circle(2.2pt);\fill[orange!85!black] (-0.7320508075688772,0) circle(2.2pt);\fill[orange!85!black] (2.732050807568877,0) circle(2.2pt);\draw[gray,dashed] (0,3)--(1,3);\node[left] at (0,3){$3$};\node[above] at (1,3){$(1,3)$};\end{scope}\node[anchor=west,align=left,text width=5.0cm] at (8.4,4.0) {$f=-(x-1)^2+3$\\[9pt]$a=-1<0,\ k=3$\\[9pt]最大值3\\[9pt]$\Delta=-4ak=12>0$\\[9pt]仍有两个实数交点};
\node[text width=12.6cm,align=center] at (7.1,.65) {先看a的符号，再判断k是最大值还是最小值，不能只看k的正负。};
\end{tikzpicture}\end{center}


；判断前先看 $a$ 的正负。


\end{trapbox}
\clearpage
\begin{trapbox}
\textbf{\textcolor{CTrap}{易错5：知道两根，就以为函数已完全确定。}}
\begin{center}\begin{tikzpicture}[font=\small,>=stealth,line join=round]
\draw[gray!40,rounded corners=3pt] (0,0) rectangle (14.2,8.3);\node at (7.1,7.700000000000001) {零点相同，a不同，曲线的高度与方向仍会不同};
\begin{scope}[shift={(1.7000000000000002,3.6214285714285714)},x=1.2cm,y=0.6571428571428571cm]\draw[->,gray] (-0.5,0)--(4.5,0) node[right]{$x$};\draw[->,gray] (0,-3)--(0,4) node[above]{$y$};\draw[blue,thick,solid,domain=0.4:3.6,samples=180] plot(\x,{(\x-1)*(\x-3)});\draw[teal,thick,solid,domain=0.4:3.6,samples=180] plot(\x,{2*(\x-1)*(\x-3)});\draw[violet,thick,solid,domain=0.4:3.6,samples=180] plot(\x,{-(\x-1)*(\x-3)});\fill[orange!85!black] (1,0) circle(2.2pt);\fill[orange!85!black] (3,0) circle(2.2pt);\fill[orange!85!black] (2,-1) circle(2.2pt);\fill[orange!85!black] (2,-2) circle(2.2pt);\fill[orange!85!black] (2,1) circle(2.2pt);\foreach \v in {1,2,3}{\node[below=4pt] at (\v,0){$\v$};}\end{scope}\node[anchor=west,align=left,text width=5.0cm] at (8.4,4.0) {蓝：$a=1$，顶点高 $-1$\\[9pt]青：$a=2$，顶点高 $-2$\\[9pt]紫：$a=-1$，顶点高 $1$\\[9pt]零点都为 $1,3$\\[9pt]还须额外条件确定a};
\node[text width=12.6cm,align=center] at (7.1,.65) {同尺度：两根式中的a控制开口，不能从两个零点中省去。};
\end{tikzpicture}\end{center}


$(x-1)(x-3)$ 与 $2(x-1)(x-3)$ 具有相同交点，但高度不同；还需要求出 $a$。


\end{trapbox}
\clearpage
\begin{trapbox}
\textbf{\textcolor{CTrap}{易错6：只知道对称轴，就以为顶点也确定。}}
\begin{center}\begin{tikzpicture}[font=\small,>=stealth,line join=round]
\draw[gray!40,rounded corners=3pt] (0,0) rectangle (14.2,8.3);\node at (7.1,7.700000000000001) {轴相同，不代表顶点和开口系数已经确定};
\begin{scope}[shift={(1.7000000000000002,2.033333333333333)},x=1.2cm,y=0.3833333333333333cm]\draw[->,gray] (-0.5,0)--(4.5,0) node[right]{$x$};\draw[->,gray] (0,-1)--(0,11) node[above]{$y$};\draw[blue,thick,solid,domain=0.6:3.4,samples=180] plot(\x,{(\x-2)^2});\draw[teal,thick,solid,domain=0.6:3.4,samples=180] plot(\x,{3*(\x-2)^2+4});\fill[orange!85!black] (2,0) circle(2.2pt);\fill[orange!85!black] (2,4) circle(2.2pt);\draw[violet,dashed] (2,-.8)--(2,10.5);\node[below right] at (2,0){$(2,0)$};\node[below right=4pt] at (2,4){$(2,4)$};\end{scope}\node[anchor=west,align=left,text width=5.0cm] at (8.4,4.0) {蓝：$(x-2)^2$\\[9pt]青：$3(x-2)^2+4$\\[9pt]轴均为 $x=2$\\[9pt]顶点 $(2,0),(2,4)$\\[9pt]a与k仍需条件};
\node[text width=12.6cm,align=center] at (7.1,.65) {同尺度：知道对称轴只确定h，顶点高度k和a仍可改变。};
\end{tikzpicture}\end{center}


若轴是 $x=2$，则 $y=(x-2)^2$ 和 $y=3(x-2)^2+4$ 都符合条件；$a,k$ 仍未知。


\end{trapbox}
\clearpage
\begin{trapbox}
\textbf{\textcolor{CTrap}{易错7：混淆与 $x$ 轴和 $y$ 轴的交点。}一般式中的 $c=f(0)$ 说明 $(0,c)$ 在 \textbf{$y$ 轴}}
\begin{center}\begin{tikzpicture}[font=\small,>=stealth,line join=round]
\draw[gray!40,rounded corners=3pt] (0,0) rectangle (14.2,8.3);\node at (7.1,7.700000000000001) {三种写法，对应同一条抛物线};
\begin{scope}[shift={(1.7000000000000002,2.9642857142857144)},x=1.2cm,y=0.6571428571428571cm]\draw[->,gray] (-0.5,0)--(4.5,0) node[right]{$x$};\draw[->,gray] (0,-2)--(0,5) node[above]{$y$};\draw[blue,thick,solid,domain=-0.2:4.2,samples=180] plot(\x,{(\x-2)^2-1});\fill[orange!85!black] (0,3) circle(2.2pt);\fill[orange!85!black] (1,0) circle(2.2pt);\fill[orange!85!black] (3,0) circle(2.2pt);\fill[orange!85!black] (2,-1) circle(2.2pt);\draw[violet,dashed] (2,-.95)--(2,4.5);\node[left] at (0,3){$(0,3)$};\node[below=5pt] at (1,0){$1$};\node[below=5pt] at (3,0){$3$};\node[below=4pt] at (2,-1){$(2,-1)$};\end{scope}\node[anchor=west,align=left,text width=5.0cm] at (8.4,4.0) {$x^2-4x+3$\\[9pt]$=(x-2)^2-1$\\[9pt]$=(x-1)(x-3)$\\[9pt]顶点 $(2,-1)$\\[9pt]交点 $(1,0),(3,0)$};
\node[text width=12.6cm,align=center] at (7.1,.65) {一般式读纵轴交点，顶点式读顶点，两根式读横轴交点。};
\end{tikzpicture}\end{center}


上；与 $x$ 轴交点必须令 $y=0$ 才能得到。


\end{trapbox}
\clearpage
\begin{trapbox}
\textbf{\textcolor{CTrap}{易错8：公式 $\Delta=-4ak$ 用错符号。}}
\begin{center}\begin{tikzpicture}[font=\small,>=stealth,line join=round]
\draw[gray!40,rounded corners=3pt] (0,0) rectangle (14.2,8.3);\node at (7.1,7.700000000000001) {顶点式中的k：开口向下时是最大值};
\begin{scope}[shift={(3.0285714285714285,3.375)},x=1.0714285714285714cm,y=0.575cm]\draw[->,gray] (-1.8,0)--(3.8,0) node[right]{$x$};\draw[->,gray] (0,-3)--(0,5) node[above]{$y$};\draw[blue,thick,solid,domain=-1.3:3.3,samples=180] plot(\x,{-(\x-1)^2+3});\fill[orange!85!black] (1,3) circle(2.2pt);\fill[orange!85!black] (-0.7320508075688772,0) circle(2.2pt);\fill[orange!85!black] (2.732050807568877,0) circle(2.2pt);\draw[gray,dashed] (0,3)--(1,3);\node[left] at (0,3){$3$};\node[above] at (1,3){$(1,3)$};\end{scope}\node[anchor=west,align=left,text width=5.0cm] at (8.4,4.0) {$f=-(x-1)^2+3$\\[9pt]$a=-1<0,\ k=3$\\[9pt]最大值3\\[9pt]$\Delta=-4ak=12>0$\\[9pt]仍有两个实数交点};
\node[text width=12.6cm,align=center] at (7.1,.65) {先看a的符号，再判断k是最大值还是最小值，不能只看k的正负。};
\end{tikzpicture}\end{center}


例如 $y=-(x-1)^2+3$，$a=-1,k=3$，$\Delta=-4(-1)\cdot3=12>0$；虽然 $k>0$，但它的顶点是最高点，不是最低点。


\end{trapbox}
\clearpage
\begin{trapbox}
\textbf{\textcolor{CTrap}{易错9：把“已知三个点”误读成“任意三个点一定唯一确定二次函数”。}}
\begin{center}\begin{tikzpicture}[font=\small,>=stealth,line join=round]
\draw[gray!40,rounded corners=3pt] (0,0) rectangle (14.2,8.3);\node at (7.1,7.700000000000001) {三个点可能只能确定直线，二次项系数为0};
\begin{scope}[shift={(2.3000000000000003,2.11)},x=1.5cm,y=0.9199999999999999cm]\draw[->,gray] (-0.8,0)--(3.2,0) node[right]{$x$};\draw[->,gray] (0,-0.5)--(0,4.5) node[above]{$y$};\draw[blue,thick,solid,domain=-0.4:2.7,samples=180] plot(\x,{\x+1});\fill[orange!85!black] (0,1) circle(2.2pt);\fill[orange!85!black] (1,2) circle(2.2pt);\fill[orange!85!black] (2,3) circle(2.2pt);\node[left] at (0,1){$(0,1)$};\node[above left] at (1,2){$(1,2)$};\node[above left] at (2,3){$(2,3)$};\end{scope}\node[anchor=west,align=left,text width=5.0cm] at (8.4,4.0) {$(0,1),(1,2),(2,3)$\\[9pt]都在 $y=x+1$ 上\\[9pt]代入 $ax^2+bx+c$\\[9pt]只能解得 $a=0$\\[9pt]不是真正的二次函数};
\node[text width=12.6cm,align=center] at (7.1,.65) {横坐标互异还不够：求出系数后，必须核对a不等于0。};
\end{tikzpicture}\end{center}


三个点的横坐标必须互不相同，且算出来的二次项系数不为 $0$。例如 $(0,1),(1,2),(2,3)$ 三点都在直线 $y=x+1$ 上，经过它们的次数不超过二的式子只有 $y=x+1$，不是真正的二次函数。


\end{trapbox}
\clearpage
\begin{trapbox}
\textbf{\textcolor{CTrap}{易错10：把 $x_1,x_2$ 的顺序当成重要条件。}}
\begin{center}\begin{tikzpicture}[font=\small,>=stealth,line join=round]
\draw[gray!40,rounded corners=3pt] (0,0) rectangle (14.2,8.3);\node at (7.1,7.700000000000001) {平方差：两个零点关于中点等距};
\begin{scope}[shift={(1.5285714285714287,3.9499999999999997)},x=0.8571428571428571cm,y=0.7666666666666666cm]\draw[->,gray] (-0.5,0)--(6.5,0) node[right]{$x$};\draw[->,gray] (0,-3)--(0,3) node[above]{$y$};\draw[blue,thick,solid,domain=0.2:5.8,samples=180] plot(\x,{.5*(\x-1)*(\x-5)});\fill[orange!85!black] (1,0) circle(2.2pt);\fill[orange!85!black] (5,0) circle(2.2pt);\fill[orange!85!black] (3,-2) circle(2.2pt);\draw[violet,dashed] (3,-2.6)--(3,2.6);\draw[<->,orange!85!black] (1,1)--(3,1) node[midway,above]{$d=2$};\draw[<->,orange!85!black] (3,1)--(5,1) node[midway,above]{$d=2$};\foreach \v/\L in {1/{x_1=1},3/{m=3},5/{x_2=5}}{\node[below=4pt] at (\v,-2.7){$\L$};}\end{scope}\node[anchor=west,align=left,text width=5.0cm] at (8.4,4.0) {图取 $a=\frac12$\\[9pt]$x_1=1,\ x_2=5$\\[9pt]$m=3,\ d=2$\\[9pt]$f=\frac12[(x-3)^2-4]$\\[9pt]顶点 $(3,-2)$};
\node[text width=12.6cm,align=center] at (7.1,.65) {两根分别是m减d与m加d；将两括号相乘，就出现以m为轴的平方项。};
\end{tikzpicture}\end{center}


交换两个括号，乘积不变。两根的和与积、对称轴中点都不会因为交换顺序而改变。


\end{trapbox}
